\documentclass[11pt,a4paper]{article}
\usepackage[T1]{fontenc}
\usepackage{lmodern}
\usepackage{microtype}
\usepackage[margin=1in]{geometry}
\usepackage{amsmath,amssymb,amsthm}
\usepackage{graphicx}
\usepackage{booktabs}
\usepackage[numbers]{natbib}
\usepackage[hidelinks]{hyperref}

\newtheorem{theorem}{Theorem}[section]
\newtheorem{lemma}[theorem]{Lemma}
\newtheorem{proposition}[theorem]{Proposition}
\newtheorem{corollary}[theorem]{Corollary}
\theoremstyle{definition}
\newtheorem{definition}[theorem]{Definition}
\newtheorem{example}[theorem]{Example}
\newtheorem{exercise}[theorem]{Exercise}
\theoremstyle{remark}
\newtheorem{remark}[theorem]{Remark}

\title{<<SCIENT_TITLE>>}
\author{<<SCIENT_AUTHOR_BLOCK>>}
\date{}

\begin{document}
\maketitle

\tableofcontents

\section{Introduction}
% Guide: Introduce the topic and the main question of the lecture.
\par

\begin{definition}
% Guide: Define the first concept.
\end{definition}

\begin{theorem}
% Guide: State the main result.
\end{theorem}

\begin{proof}
% Guide: Prove it.
\end{proof}

\begin{example}
% Guide: Work through an example.
\end{example}

\begin{exercise}
% Guide: Set an exercise for the reader.
\end{exercise}

\bibliographystyle{plainnat}
\bibliography{references}

\end{document}
